\section*{Chapitre \ref{ch:g9:solids} --- Solides et volumes}

\begin{solution}{exo:g9:solids:1}
\omterm{def:g5:solids:def}{Pavé}~: $V = 4 \times 5 \times 12 = 240$.

\omterm{def:g7:areas:prism}{Cylindre}~: $V = \pi r^2 h = \pi \times 9 \times 7 = 63\pi \approx 198$.
\end{solution}

\begin{solution}{exo:g9:solids:2}
\omterm{def:g8:solids:def}{Cône}~: $V = \frac13 \pi r^2 h = \frac13 \pi \times 25 \times 12 =
100\pi \approx 314$.

\omterm{def:g8:solids:def}{Pyramide}~: \omterm{def:g6:measure:area}{aire} de base $B = 8 \times 3 = 24$, donc
$V = \frac13 \times 24 \times 10 = 80$.
\end{solution}

\begin{solution}{exo:g9:solids:3}
\omterm{def:g6:measure:volume}{Volume}~: $V = \frac43 \pi \times 6^3 = \frac43 \pi \times 216 = 288\pi$.

\omterm{def:g6:measure:area}{Aire} de surface~: $4\pi \times 6^2 = 144\pi$.
\end{solution}

\begin{solution}{exo:g9:solids:4}
\omterm{def:g3:shapes:shapes}{Rayon} de la \omterm{prop:g9:solids:sections}{section}~: $\sqrt{r^2 - d^2} = \sqrt{100 - 64} = \sqrt{36}
= 6$.
\end{solution}

\begin{solution}{exo:g9:solids:5}
\omterm{def:g9:thales:scaling}{Facteur d'échelle} du \omterm{def:g8:solids:def}{cône} entier vers le petit~:
$k = \frac{9}{12} = \frac34$. \omterm{def:g3:shapes:shapes}{Rayon} de \omterm{prop:g9:solids:sections}{section}~: $8 \times \frac34 = 6$.

\omterm{def:g6:measure:volume}{Volume} du \omterm{def:g8:solids:def}{cône} entier~: $\frac13\pi \times 64 \times 12 = 256\pi$.
\omterm{def:g6:measure:volume}{Volume} du petit \omterm{def:g8:solids:def}{cône}~: $256\pi \times \left(\frac34\right)^3 =
256\pi \times \frac{27}{64} = 108\pi \approx 339$.
\end{solution}

\begin{solution}{exo:g9:solids:6}
\omterm{def:g6:measure:volume}{Volume} de la bille~: $\frac43\pi \times 2^3 = \frac{32\pi}{3}$ cm$^3$.
Ce \omterm{def:g6:measure:volume}{volume} se répartit sur l'\omterm{def:g6:measure:area}{aire} de base $\pi \times 3^2 = 9\pi$ cm$^2$,
donc le niveau monte de
\[
\frac{32\pi/3}{9\pi} = \frac{32}{27} \approx 1.2 \text{ cm}
\]
(environ $12$ mm).
\end{solution}

\begin{solution}{exo:g9:solids:7}
Le petit moule est la réduction du grand par $k = \frac36 =
\frac12$, donc son \omterm{def:g6:measure:volume}{volume} est $\left(\frac12\right)^3 = \frac18$ du
grand~: la recette remplit $8$ petits moules.
\end{solution}

\begin{solution}{exo:g9:solids:8}
$V = \frac13 \times 230^2 \times 147 = \frac13 \times 52\,900 \times 147
= 52\,900 \times 49 = 2\,592\,100$ m$^3$ --- environ $2.6$ millions de
mètres cubes.
\end{solution}

\begin{solution}{exo:g9:solids:9}
\emph{1.} L'eau forme un \omterm{def:g8:solids:def}{cône} réduit par $k = \frac12$ (\omterm{def:g5:solids:def}{sommet} en bas,
\omterm{def:g3:division:half}{moitié} de la hauteur), donc son \omterm{def:g6:measure:volume}{volume} est $\left(\frac12\right)^3 =
\frac18$ de celui de l'entonnoir~: un huitième, bien moins que la
\omterm{def:g3:division:half}{moitié}~!

\emph{2.} L'eau jusqu'à la hauteur $d$ forme un \omterm{def:g8:solids:def}{cône} réduit par
$k = \frac{d}{12}$, de \omterm{def:g6:measure:volume}{volume} $k^3$ fois celui de l'entonnoir. La \omterm{def:g3:division:half}{moitié}
du \omterm{def:g6:measure:volume}{volume} signifie $k^3 = \frac12$, c'est-à-dire
$\left(\frac{d}{12}\right)^3 = \frac12$. Puis
$\frac{d}{12} = \sqrt[3]{0.5} \approx 0.7937$, donc
$d \approx 9.5$ cm~: la seconde \omterm{def:g3:division:half}{moitié} du \omterm{def:g6:measure:volume}{volume} n'occupe que les $2.5$
cm du haut --- la partie large du \omterm{def:g8:solids:def}{cône} contient la plupart de l'eau.
\end{solution}

\begin{solution}{pb:g9:solids:1}
\textbf{1.} $V_{\text{cyl}} = \pi r^2 \times 2r = 2\pi r^3$.

\textbf{2.} $\dfrac{V_{\text{sphère}}}{V_{\text{cyl}}}
= \dfrac{\frac43 \pi r^3}{2 \pi r^3} = \dfrac23$. Le $\pi$ et le
$r^3$ se simplifient~: la proportion est \emph{universelle} ---
vraie pour une bille comme pour une planète. Une relation entre
formes, et non entre nombres~: exactement le genre de vérité qui
mérite d'être gravée dans la pierre.

\textbf{3.} Sphère~: $4\pi r^2$. \omterm{def:g7:areas:prism}{Cylindre}~: latérale
$2\pi r \times 2r = 4\pi r^2$, plus les deux disques
$2 \times \pi r^2$~: total $6\pi r^2$. Rapport~:
$\frac{4\pi r^2}{6\pi r^2} = \frac23$ --- les mêmes deux tiers,
pour les surfaces. (Et une prime~: l'\omterm{def:g6:measure:area}{aire} de la sphère est
exactement égale à l'\omterm{def:g6:measure:area}{aire} latérale du \omterm{def:g7:areas:prism}{cylindre}.)

\textbf{4.} \omterm{def:g3:division:remainder}{Reste}~: $2\pi r^3 - \frac43\pi r^3 = \frac23 \pi r^3$.
Deux \omterm{def:g8:solids:def}{cônes} de \omterm{def:g3:shapes:shapes}{rayon} $r$ et de hauteur $r$~:
$2 \times \frac13 \pi r^2 \times r = \frac23 \pi r^3$. Égalité.

\textbf{5.} Ballon~: $\frac43 \pi \times 12^3 \approx 7\,200$
cm$^3$~; boîte~: $2\pi \times 12^3 \approx 10\,900$ cm$^3$. Vide~:
le tiers de la boîte, environ $33\,\%$ --- garanti par la question
2, quelle que soit la taille du ballon.

\textbf{6.} Une demi-sphère~:
$\frac12 \times \frac43 \pi \times 10^3 = \frac23 \pi \times
1\,000 \approx 2\,094$ cm$^3 \approx 2.1$ L.

\textbf{7.} \omterm{def:g8:solids:def}{Cône}~: $\frac13 \pi \times 10^2 \times 10 =
\frac{1\,000\pi}{3} \approx 1\,047$ cm$^3$. Demi-boule~:
$\frac{2\,000\pi}{3} \approx 2\,094$ cm$^3$. \omterm{def:g7:areas:prism}{Cylindre}~:
$1\,000\pi \approx 3\,142$ cm$^3$. Rapports~:
$\frac{1000\pi}{3} : \frac{2000\pi}{3} : \frac{3000\pi}{3}
= 1 : 2 : 3$ exactement.

\textbf{8.} $\frac43 \pi \times 27 = \pi \times 9 \times h$ donne
$h = \frac43 \times 3 = 4$ cm exactement.

\textbf{9.} \omterm{def:g3:shapes:shapes}{Rayons}~: $\frac{6\,371}{1\,737} \approx 3.67$. Les
\omterm{def:g6:measure:volume}{volumes} suivent le \omterm{def:g5:solids:def}{cube} (\cref{thm:g9:solids:scaling})~:
$3.67^3 \approx 49$. Une cinquantaine de Lunes tiennent dans la
Terre.

\textbf{10.} Les surfaces suivent le carré~:
$3.67^2 \approx 13.5$. Doubler le \omterm{def:g3:shapes:shapes}{rayon} d'un ballon multiplie son
\omterm{def:g6:measure:volume}{volume} par $2^3 = 8$ mais sa surface par $2^2 = 4$ seulement~: à
mesure que les choses grandissent, le \omterm{def:g6:measure:volume}{volume} (le poids, le contenu,
la chaleur produite) distance la surface (la peau, la matière, la
chaleur perdue) --- c'est la loi du \omterm{def:g5:solids:def}{carré-cube}.

\textbf{11.} Poids~: $\times 10^3 = 1\,000$. \omterm{prop:g9:solids:sections}{Section} des os~:
$\times 10^2 = 100$. Pression --- le poids par unité d'\omterm{def:g6:measure:area}{aire} d'os~:
$\times \frac{1000}{100} = 10$. Des os conçus pour une pression
humaine en reçoivent dix fois plus~: le squelette du géant cède sous
son propre poids. Les géants sont géométriquement impossibles~; les
grands animaux ont besoin d'os démesurément épais (comparer les
pattes d'un éléphant à celles d'une gazelle).

\textbf{12.} Poids~: $\div 100^3 = 10^6$. Force (\omterm{prop:g9:solids:sections}{section}
musculaire)~: $\div 100^2 = 10^4$. Rapport force sur poids~:
$\times \frac{10^6}{10^4} = 100$. La fourmi qui soulève cinquante
fois son poids n'est pas une super-athlète --- elle est simplement
petite~; un être humain réduit à la taille d'une fourmi en ferait
autant.

\textbf{13.} $\dfrac{S}{V} = \dfrac{4\pi r^2}{\frac43\pi r^3}
= \dfrac{3}{r}$~: pour $r = 1$ le rapport vaut $3$, pour $r = 2$ il
vaut $1.5$ --- diviser la taille par deux double la surface
disponible par unité de \omterm{def:g6:measure:volume}{volume} producteur de chaleur. Les petits
corps perdent leur chaleur~: la souris doit manger sans cesse,
l'ours peut hiberner, et le bébé --- petite sphère, grand rapport
$\frac Sr$ --- a besoin de son bonnet.

\textbf{14.} $\left(\frac54\right)^3 = \frac{125}{64} \approx
1.95$~: la grosse orange contient près de \emph{deux fois} plus de
fruit pour le même prix. Il faut acheter du \omterm{def:g3:shapes:shapes}{rayon}.

\textbf{15.} La gravure dit ceci~: la sphère vaut $\frac23$ du
\omterm{def:g7:areas:prism}{cylindre} en \omterm{def:g6:measure:volume}{volume}, \emph{et} $\frac23$ de sa surface totale ---
deux proportions exactes et universelles reliant le plus rond des
\omterm{def:g5:solids:def}{solides} au plus simple. La réponse à Cicéron~: parce qu'un théorème
est le seul monument que ni les armées ni les siècles n'érodent ---
les machines d'Archimède ont brûlé avec Syracuse, mais les deux
tiers valent toujours exactement deux tiers.
\end{solution}
